\section*{Bab \ref{ch:g4:classifying-animals} --- Mengklasifikasikan Hewan}

\begin{solution}{exo:g4:classifying-animals:1}
Seekor \omterm{def:g4:classifying-animals:vertebrate}{vertebrata} punya \omterm{def:g3:bones-and-muscles:skeleton}{rangka} dalam yang dibangun mengelilingi \omterm{def:g4:classifying-animals:vertebrate}{tulang
belakang} (\omterm{prop:g3:sorting-living-things:groups}{ikan} trout, kucing); seekor \omterm{def:g4:classifying-animals:vertebrate}{invertebrata} tidak punya \omterm{def:g4:classifying-animals:vertebrate}{tulang
belakang} (siput, lebah, cacing tanah, ubur-ubur).
\end{solution}

\begin{solution}{exo:g4:classifying-animals:2}
\omterm{prop:g3:sorting-living-things:groups}{Ikan}: bersisik, bersirip, bernapas di dalam air. \omterm{prop:g3:sorting-living-things:groups}{Amfibi}: berkulit
telanjang yang lembap, \omterm{def:g2:animals-grow-up:born}{telurnya} di air, masa kecilnya sebagai \omterm{ex:g2:animals-grow-up:frog}{berudu}.
\omterm{prop:g4:classifying-animals:classes}{Reptil}: berkulit bersisik yang kering, kebanyakan mengeluarkan \omterm{def:g2:animals-grow-up:born}{telur}
bercangkang di darat. \omterm{prop:g3:sorting-living-things:groups}{Burung}: berbulu \omterm{prop:g3:sorting-living-things:groups}{burung}, berparuh, bertelur
bercangkang. \omterm{prop:g3:sorting-living-things:groups}{Mamalia}: berbulu atau berambut, anaknya diberi \omterm{def:g2:animals-grow-up:born}{susu}.
\end{solution}

\begin{solution}{exo:g4:classifying-animals:3}
Sarden: \omterm{prop:g3:sorting-living-things:groups}{ikan}. Buaya: \omterm{prop:g4:classifying-animals:classes}{reptil}. Salamander: \omterm{prop:g3:sorting-living-things:groups}{amfibi}. \omterm{prop:g3:sorting-living-things:groups}{Burung} unta: \omterm{prop:g3:sorting-living-things:groups}{burung}.
Paus: \omterm{prop:g3:sorting-living-things:groups}{mamalia}. Cacing tanah: \omterm{def:g4:classifying-animals:vertebrate}{invertebrata} --- tanpa \omterm{def:g4:classifying-animals:vertebrate}{tulang belakang},
jadi sama sekali tidak masuk \omterm{prop:g4:classifying-animals:classes}{kelas vertebrata}.
\end{solution}

\begin{solution}{exo:g4:classifying-animals:4}
Kucingnya, di dalam \omterm{prop:g3:sorting-living-things:groups}{mamalia}, di dalam \omterm{def:g4:classifying-animals:vertebrate}{vertebrata}, di dalam \omterm{def:g1:animals-around-us:animal}{hewan}.
\end{solution}

\begin{solution}{exo:g4:classifying-animals:5}
\omterm{def:g1:the-five-senses:sense}{Kulitnya} --- kering dan bersisik pada kadal, lembap dan telanjang pada
salamander --- dan \omterm{def:g2:animals-grow-up:born}{telurnya}: bercangkang dan diletakkan di darat pada
kadal, telanjang dan diletakkan di air pada salamander.
\end{solution}

\begin{solution}{exo:g4:classifying-animals:6}
\omterm{def:g4:classifying-animals:vertebrate}{Invertebrata}. Misalnya: \omterm{prop:g3:sorting-living-things:groups}{serangga}, laba-laba, siput, cacing (dan
ubur-ubur).
\end{solution}

\begin{solution}{exo:g4:classifying-animals:7}
Penguin: berbulu \omterm{prop:g3:sorting-living-things:groups}{burung}? Ya --- seekor \omterm{prop:g3:sorting-living-things:groups}{burung}; berhenti pada
pertanyaan pertama. Lumba-lumba: berbulu \omterm{prop:g3:sorting-living-things:groups}{burung}? Tidak. Berbulu atau
bersusu? Ia memberi \omterm{def:g2:animals-grow-up:born}{susu} kepada anaknya --- seekor \omterm{prop:g3:sorting-living-things:groups}{mamalia}; berhenti
di situ.
\end{solution}

\begin{solution}{exo:g4:classifying-animals:8}
Misalnya: \omterm{def:g1:animals-around-us:covering}{bulu burung}, sebuah paruh, sayap, \omterm{def:g2:animals-grow-up:born}{telur} bercangkang --- dan
sebuah \omterm{def:g4:classifying-animals:vertebrate}{tulang belakang}, sebab \omterm{prop:g3:sorting-living-things:groups}{burung} termasuk \omterm{def:g4:classifying-animals:vertebrate}{vertebrata}.
\end{solution}

\begin{solution}{exo:g4:classifying-animals:9}
Pertanyaannya berhenti pada ``\omterm{def:g2:animals-grow-up:born}{susu} bagi anaknya'': pausnya menyusui
anaknya, jadi ia seekor \omterm{prop:g3:sorting-living-things:groups}{mamalia}, hampir tak berambut atau tidak.
``Hidup di laut'' adalah \omterm{def:g2:where-animals-live:habitat}{habitat}, bukan ciri tubuh --- dan \omterm{def:g4:classifying-animals:classification}{klasifikasi}
dibangun atas apa yang \emph{dimiliki} seekor \omterm{def:g1:animals-around-us:animal}{hewan} justru karena
kebiasaan dan rumah dapat dibagi oleh yang bukan kerabat: hiu dan
lumba-lumba berbagi alamat, bukan kelas.
\end{solution}

\begin{solution}{exo:g4:classifying-animals:10}
\omterm{def:g4:classifying-animals:vertebrate}{Tulang belakangnya}. Ular punya \omterm{def:g4:classifying-animals:vertebrate}{tulang belakang} yang terdiri atas
ratusan \omterm{def:g3:bones-and-muscles:skeleton}{tulang} kecil --- seekor \omterm{def:g4:classifying-animals:vertebrate}{vertebrata}, dengan \omterm{def:g1:the-five-senses:sense}{kulit} bersisik yang
kering milik \omterm{prop:g4:classifying-animals:classes}{reptil}; cacing tanah sama sekali tidak punya \omterm{def:g4:classifying-animals:vertebrate}{tulang
belakang} --- seekor \omterm{def:g4:classifying-animals:vertebrate}{invertebrata}. Kaki yang tidak ada hanyalah ciri
yang tidak ada; pertanyaan \omterm{def:g4:classifying-animals:vertebrate}{tulang belakang} diajukan lebih dahulu.
\end{solution}

\begin{solution}{exo:g4:classifying-animals:11}
\omterm{prop:g3:sorting-living-things:groups}{Mamalia} yang menang: \omterm{def:g1:animals-around-us:covering}{bulu} dan \omterm{def:g2:animals-grow-up:born}{susu} bagi anaknya adalah berkas ciri
\omterm{prop:g3:sorting-living-things:groups}{mamalia}, dan platipus punya keduanya --- \omterm{def:g2:animals-grow-up:born}{telurnya} adalah satu ciri
yang keluar dari barisan. \omterm{def:g4:classifying-animals:classification}{Klasifikasi} menimbang seluruh berkasnya
alih-alih membiarkan satu ciri memutuskan; \omterm{def:g1:animals-around-us:animal}{hewan} yang ganjil dalam
satu hal tetap duduk di tempat yang ditunjuk sebagian besar cirinya.
(Platipus cukup ganjil sehingga \omterm{prop:g3:sorting-living-things:groups}{mamalia} yang bertelur membentuk
sub-kelompok kecilnya sendiri.)
\end{solution}
